\section{Introduction}

Beginning in July 1968, John B.~Calhoun followed a population of house mice in a closed enclosure that he called
\emph{Universe~25} \cite{Calhoun1973}.
Food, water and nest space were in excess, predators and epidemic disease were excluded, and emigration was
impossible.
From four founding pairs the colony grew roughly exponentially.
It stopped growing at about 2200 animals, well below the 3840 the nest boxes could shelter.
The last young to survive weaning were born shortly after the peak, and the population then aged out without ever
recovering.
Along the way Calhoun described a breakdown of social behaviour:
\begin{itemize}
\item withdrawn males gathered in pools at the centre of the floor;
\item nursing females attacked and abandoned their young;
\item there appeared the unwounded, non-reproducing ``beautiful ones''.
\end{itemize}
Small groups moved to uncrowded enclosures during the decline did not reproduce, even with partners raised without
crowding (Marsden's work, reported in \cite{Calhoun1973}; see also \cite{Marsden1972}).
Calhoun read the collapse as a failure of social roles and of the maturation of complex behaviour, not of resources:
``density per se was not the major factor'' \cite{Calhoun1973}.
Section~\ref{sec:observations} lists the reported facts with page references.

The interpretation of the experiment has always been contested.
Its popular message, ``density causes pathology'', departed from Calhoun's own emphasis on social roles and design
\cite{RamsdenAdams2009,Ramsden2009,AdamsRamsden2024}.
Studies of human crowding found weak or no effects of density per se \cite{Galle1972,Freedman1975}.
For modelling, the decisive limitation is that Universe~25 was a single realization.
Calhoun himself called the study ``observation and reconstruction of a process; it was historical''
\cite{Calhoun1973}.
A model tuned to the published curve therefore risks fitting one sample path of a strongly stochastic process
\cite{GwizdallaDus2026}.

We ask instead:
\begin{quote}
Which \emph{local} interaction rules, with unlimited resources and no rule that depends on the total population,
generate robustly across stochastic realizations the part of Calhoun's sequence that runs from rapid growth to
extinction?
That part includes deceleration with social deterioration, cessation of reproduction near the peak, a decline
without recovery, and differentiated classes of deteriorated animals.
\end{quote}
The latency phase and the growth from four founding pairs are part of Calhoun's sequence, but we do not claim them:
our founder protocol reproduces them only partially (Section~\ref{sec:discussion}).
Our core model is a lattice model in which each mouse carries an age, a sex and a continuous social state $s_i\in[0,1]$.
The social state deteriorates under local crowding, scales conception and pup survival, and is partly inherited.
Movement is biased towards acquaintances.
We then add a few local rules, each with a neutral value that switches it off and recovers the core model, and each motivated by an
observation of the original report (Section~\ref{sec:model-rules}). Most have one principal parameter; the refuge
rule R7 has five and was designed after seeing the first candidate (Section~\ref{sec:limitations}):
\begin{itemize}
\item a juvenile window for social plasticity (R1);
\item social learning of competence (R2);
\item crowd aversion of deteriorated animals (AV);
\item refuges of limited capacity preferred by deteriorated animals (R7).
\end{itemize}
Rule names follow the simulation code; rules tested and discarded are listed in Section~\ref{sec:model-rules}.

Our contributions are the following. We separate throughout what is preregistered from what is post hoc, and what
the rules \emph{produce} from what they \emph{encode}.
\begin{enumerate}
\item We translate Calhoun's report, checked against the original, into seventeen falsifiable patterns, five
  refinements and four anti-patterns. Each has an operational criterion evaluated run by run, following
  pattern-oriented modelling \cite{Grimm2005}. We report the discriminating power of every essential criterion
  against four null models (no crowding damage; full inheritance without social learning; the long-lived demography
  without new rules; and the calibrated candidate without its four rules) and two partial ablations (no attraction,
  and the candidate without refuges) (Table~\ref{tab:res-null}). The preregistered forms of three criteria (O4, O10, O11) pass in null models; the primary evaluator,
  not preregistered, redefines them. O4 and O10 remain generic. O11 discriminates when the late cohorts are scored at six weeks
  of age, but not when they are scored at the end of the plasticity window. An additional
  desirable pattern (O10b, not preregistered) asks for a growth phase with socially competent adults.
\item We show that the core model without the rules, with a short life span (the short-lived control), is not
  Calhoun-like.
  About half of its colonies die in the first collapse, and the others oscillate with recovering natality.
  The mean-field analysis of that model traces the oscillations to spatial pockets of healthy individuals and to
  unconditional recovery.
  Removing the plasticity window from the calibrated model brings back recovery and recurrent natality at fast
  learning rates.
\item We separate what the rules add from what the demography, the crowding damage and the design of the rules
  already imply.
  \begin{itemize}
  \item \emph{The generic envelope.} A long-lived demography alone already drives every colony to extinction, so
    extinction is \emph{not} a result of the rules. With crowding damage and the calibrated attraction, the
    candidate without its four rules also has a low peak with free space (it passes the primary O4 in 200 of 200
    runs, with a peak at 0.45 of the damage-threshold capacity and 74\% of the cells empty), an end of natality and
    no rebound. Scored at the end of the plasticity window, its late-born cohorts also fail to acquire social
    competence (O11 in 0.755 of the runs, and 0.81 in the long-lived control), through accumulated crowding damage.
    These patterns verify that the model has a boom and its end; they are not evidence for the rules.
  \item \emph{Consistency checks of the design.} Some patterns follow from the rules by construction, and we report
    them as checks of the implementation, not as findings. Transplanted late-phase animals never found a colony
    (O12), because under the plasticity window R1 an adult cannot recover and conception is proportional to the
    social state of both partners, which is essentially zero in these animals. The isolated class is persistent
    (O13p) because the refuge capacity equals the isolation threshold. There are no social deaths (A4) and the decline
    is senescent (O8) because social mortality is switched off.
  \item \emph{What the rules add and the null models lack.} Two patterns: the coexistence of crowded and
    isolated deteriorated classes (O13), which needs crowd aversion or refuges and is a calibrated property of a
    nearly sessile regime; and the persistence of the isolated class (O13p), which is a design property. Both
    discriminate against three of the four null models, not against full inheritance without social learning. The
    failure of the late cohorts (O11) is not counted here. Scored at six weeks of age, as in the primary
    evaluator, it discriminates against all four null models, but that is partly definitional: pups are born
    below the deterioration threshold and are still inside the plasticity window. Scored at the end of the window,
    \textsf{C1} still passes (0.995), the dependence on the birth state disappears, and O11 discriminates against
    neither the long-lived nor the rule-free null model (0.81 and 0.755), whose late cohorts end even lower than
    those of \textsf{C1} (median $\bar s$ 0.07 and 0.06 against 0.13). The failure of the late cohorts is therefore mostly crowding damage.
\item \emph{What the dynamics is.} A boom of one or two generations, followed by a plateau without births or
    deaths and extinction by senescence. About 70\% of the births are to the founding animals, and the net
    reproductive number falls from 3.8 to 0.4 in a single generation (reference ensemble, 200 runs,
    Appendix~\ref{app:calib-generations}): pups are born below the deterioration threshold and learn from adults
    who have already deteriorated. The phases appear in Calhoun's order, but B and C are compressed, and we do not
    claim a cascade over many generations. The ratio of the deterioration time to the generation time, measured
    on animals past the plasticity window so that newborns do not dilute it, is $\Gamma=t_{\rm det}/T_{\rm gen}$
    $\approx1.1$ (1.2 on the founders), against about 3.5 in Universe~25 with a different operational definition; it is a heuristic
    comparison of time scales, not the cause of the shallow boom. With
    initial ages spread over 4--52 or 4--80 weeks the classes, the low peak and the failure of the late cohorts
    survive, but the timing of the end of natality relative to the peak (O5) does not (item~6).
  \end{itemize}
\item We identify which rules are needed, conditionally on the evaluation thresholds.
  With the class threshold of 10\%, isolation at $m\le2$, persistence of four weeks and the ensemble threshold of 0.8,
  the minimal set for the primary outcome is R1, R2 and R7 (candidate \textsf{C1}$'$, primary outcome in
  0.835 [0.78, 0.88] of 200 runs, borderline because the lower Wilson bound is below 0.8). With a class threshold of
  15\% it falls to 0.33, so minimality is not a property of the rules alone. Crowd aversion (AV) is a contributing
  rule: adding it (candidate \textsf{C1}) makes the primary outcome more reliable, through the generic criterion
  O5, and the decline senescent (O8); it \emph{reduces} the free space at the peak (16\% of the cells empty with
  AV, 52\% without). The contrast between \textsf{C1} and \textsf{C1}$'$ was not preregistered.
\item We test seven preregistered hypotheses (plan fixed on 1 October 2026, before the production runs; \texttt{docs/preregistered\_plan.md}).
  With the primary criteria the null hypothesis is rejected for six after Holm correction and the quantitative
  prediction holds for four (Q1, Q3, Q4, Q7); with the preregistered criteria the counts are five and four, and Q2 is
  not rejected. The only change that produces the rejection of Q2 is the post hoc choice of the end of natality in
  O5 (the 99th percentile of the births instead of the last birth), made after exploratory runs in which the
  alternative lowered the primary outcome of \textsf{C1} to about one half (Table~\ref{tab:deviations}). Two of the
  four predictions that hold are consistency checks, Q1 of R1 (above) and Q4 of the demography, and the
  preregistered contrast of Q3 is confounded by the demography. Q7 is confirmed only in the restricted sense stated
  in Section~\ref{sec:res-Q}. Removing crowd aversion lowers the primary outcome significantly with the primary criteria, but by much less than predicted; the predicted dependence on the length of the plasticity window (Q6)
  is not supported, because with the primary O11 the window has little or no effect on the boundary; and a
  predicted coexistence band (Q5) is refuted. Every deviation from the preregistered plan is listed in
  Table~\ref{tab:deviations}.
\item We report robustness in four senses. Against stochasticity it is high. Locally around \textsf{C1} it holds
  in most of the phase planes, with a nest period for pups (0.835 [0.78, 0.88], borderline) and with an
  asynchronous update (1.00). Globally it is limited and relative to the sampled box: the mean primary outcome over
  two 15-factor Latin hypercubes is 0.040 [0.026, 0.059], with 2.0\% of the points Calhoun-like, and the Sobol'
  indices are conditional on the factors fixed after screening. Demographically the primary outcome falls to 0.55
  with a background mortality of 0.003 per week and to 0.63 with initial ages spread over 4--80 weeks. Natality
  still ends with the population near its maximum, but the plateau starts earlier, so the last births fall more
  than half a peak time after its start and O5 becomes marginal (in 83 and 73 of the 90 and 74 runs that fail; the
  others fail it by the same timing clause). O5 is generic, and the patterns that the rules add (O13, O13p), as well as
  O11, pass in every one of these runs (Section~\ref{sec:limitations}).
\end{enumerate}
The model is documented with the ODD protocol \cite{Grimm2020} (Appendix~\ref{app:odd}), and the code is available in the repository \repourl.

\section{Related work}

\paragraph{Calhoun's experiments and their reception.}
Calhoun's programme of crowding experiments began with wild Norway rats in an outdoor enclosure.
There the population stabilized at about 150 adults, far below what adult survival would have allowed, because
of very high infant mortality \cite{Calhoun1962,Calhoun1963}.
It continued with laboratory ``universes'' for rats and mice, of which Universe~25 is the best known
\cite{Calhoun1973}.
His theoretical writings stress optimal group size and the saturation of social roles rather than physical density
\cite{Calhoun1971}.
The scientific and cultural history of these experiments, including the gap between Calhoun's interpretation and its
popular uptake, is treated in detail in \cite{RamsdenAdams2009,Ramsden2009,AdamsRamsden2024,Dugatkin2024}.
Earlier confined-population studies of house mice and other rodents had already reported that, with abundant food,
population growth stops through reduced litter survival and reproduction mediated by social interference
\cite{Southwick1955,Lidicker1965}.
This is the empirical basis of the coupling between social state and newborn survival in our model.

\paragraph{Models of Universe~25.}
Quantitative modelling of the experiment is very recent.
Gwizdałła and Duś \cite{GwizdallaDus2026} proposed an agent-based simulation with a daily time step and
empirically motivated life tables and litter-size distributions.
They calibrate the growth phase against the reported population and number of adults at day 315.
The decline is produced by deviation factors that are triggered when the total population exceeds critical values,
together with prescribed probabilities of litter abandonment and of becoming a ``beautiful one''.
They obtain the qualitative shape of the curve and approximate the timing of the peak.
They also emphasize the large dispersion across replicates and the lack of a methodological basis for the choice of
rules.
Mantzaris \cite{Mantzaris2025} attributes the collapse to the progressive loss of uncertainty about the social
hierarchy in a fully visible enclosure, implemented through a game-theoretic payoff matrix whose parameters change
across generations.
Roman \cite{Roman2025} fits maximum-entropy unimodal curves to the population time series without a mechanistic model.

Our approach differs from these in three respects:
\begin{enumerate}
\item all rules are local in space and use no global population thresholds;
\item the deterioration and its maternal transmission are dynamical variables rather than prescribed
  probabilities;
\item validation is based on multiple qualitative patterns evaluated as distributions over stochastic replicates,
  not on curve fitting.
\end{enumerate}

\paragraph{Social stress, maternal effects and population cycles.}
The idea that crowding acts through physiological stress on reproduction goes back to Christian's
adreno-pituitary hypothesis \cite{Christian1950}.
In the population-cycle literature, the transmission of individual quality from mother to offspring---the
\emph{maternal effect hypothesis}---produces delayed density dependence and cycles in simple two-variable models of
density and mean quality \cite{GinzburgTaneyhill1994,InchaustiGinzburg1998}.
Field evidence shows that maternal stress reduces reproduction and is carried over to the offspring
\cite{Sheriff2009}.
Our inherited social state is a spatially explicit, individual-based version of this mechanism.
Consistent with that theory, its mean-field limit has an oscillatory regime, which we show is not the one realized
in Universe~25.

\paragraph{Allee effects and stochastic extinction.}
Because reproduction requires that partners meet in the same cell, the model has a component Allee effect at low
density \cite{StephensSutherland1999,Courchamp2008}.
This has to be distinguished from social deterioration as an explanation for the absence of recovery.
Extinction is ultimately driven by demographic stochasticity acting on an absorbing state.
Mean times to extinction and their scaling with population size, quasi-stationary behaviour and extinction from the
troughs of cycles are reviewed in \cite{Ovaskainen2010}.
These provide the reference for our survival analysis of extinction times.

\paragraph{Methodology for agent-based models.}
We describe the model with the ODD protocol \cite{Grimm2020} and validate it in the spirit of pattern-oriented
modelling \cite{Grimm2005}.
For sensitivity analysis of stochastic agent-based models we follow the recommendations of ten Broeke et al.\
\cite{tenBroeke2016}: exploratory one-factor-at-a-time analysis to uncover mechanisms, followed by variance-based
global indices \cite{Saltelli2010}.
The number of replicates is chosen by the stability of output statistics \cite{Lorscheid2012}.
Although our aim is qualitative, the region of parameter space compatible with the reported patterns can be
characterized with approximate Bayesian computation using pattern-based summary statistics
\cite{Hartig2011,vanderVaart2015}.
